\documentclass[11pt]{article}
\usepackage{explorer}
\newcommand{\ExplorerLanguage}{FR}
\newcommand{\ExplorerNumber}{TRINÔME}
\newcommand{\ExplorerKind}{Présenter et explorer}
\titlespacing*{\section}{0pt}{8pt}{4pt}
\setlength{\parskip}{5pt}\linespread{1.06}
\begin{document}
\ExplorerTitle{Première spécialité}{Trinôme du second degré}
\ExplorerSection{Forme développée}
Pour $f(x)=ax^2+bx+c$, avec $a,b,c\in\mathbb R$ et $a\ne0$, le domaine est $\mathbb R$. Le coefficient $a$ détermine le sens d’ouverture de la parabole.\par
\ExplorerSection{Forme canonique}
$f(x)=a(x-\alpha)^2+\beta$, où $\alpha=-\dfrac{b}{2a}$ et $\beta=f(\alpha)=c-\dfrac{b^2}{4a}$. Le sommet est $S(\alpha,\beta)$ et l’axe de symétrie est $x=\alpha$. Si $a>0$, $f$ décroît sur $]-\infty;\alpha]$ et croît sur $[\alpha;+\infty[$ ; $\beta$ est le minimum. Si $a<0$, les variations sont inversées et $\beta$ est le maximum.\par
\ExplorerSection{Discriminant et racines}
$\Delta=b^2-4ac$ et $\beta=-\dfrac{\Delta}{4a}$. Si $\Delta>0$, les deux racines sont $\dfrac{-b-\sqrt\Delta}{2a}$ et $\dfrac{-b+\sqrt\Delta}{2a}$ : les ordonner avant de les appeler $x_1<x_2$ (l’ordre des deux formules dépend du signe de $a$). Si $\Delta=0$, la racine double est $x_0=\alpha$. Si $\Delta<0$, il n’y a aucune racine réelle.\par
\ExplorerSection{Factorisation et signe}
Si $\Delta>0$, $f(x)=a(x-x_1)(x-x_2)$. À l’extérieur des racines, les deux facteurs sont de même signe ; entre elles, ils sont de signes opposés. Le trinôme est donc du signe de $a$ à l’extérieur, du signe opposé entre les racines, et nul aux racines. Si $\Delta=0$, $f(x)=a(x-x_0)^2$ est du signe de $a$ sauf en $x_0$. Si $\Delta<0$, $f$ est strictement du signe de $a$ sur $\mathbb R$ et ne se factorise pas en facteurs du premier degré à coefficients réels.\par
\ExplorerSection{Exemple expliqué}
$2x^2-4x-6=2(x-1)^2-8=2(x+1)(x-3)$. Le sommet est $(1,-8)$, $\Delta=64$ et les racines sont $-1$ et $3$. Ainsi $f(x)\le0$ équivaut à $x\in[-1;3]$.\par
\begin{center}\begin{tikzpicture}[x=.65cm,y=.13cm]\draw[->,gray](-3,0)--(5,0)node[right]{$x$};\draw[->,gray](0,-10)--(0,12)node[above]{$y$};\draw[violet,thick,domain=-2.15:4.15,samples=100]plot(\x,{2*(\x-1)^2-8});\draw[mint,dashed](1,-9)--(1,11);\fill[mint](-1,0)circle(2pt)node[below]{$-1$};\fill[mint](3,0)circle(2pt)node[below]{$3$};\fill[violet](1,-8)circle(2pt)node[right]{$S(1,-8)$};\end{tikzpicture}\end{center}
\ExplorerSection{Chercher et justifier}
Étudier les variations, les racines et le signe de $-x^2+2x+3$, de $x^2-2$ et de $x^2+x+1$. Donner les racines sous forme exacte, puis une approximation clairement indiquée par $\approx$. Vérifier les résultats en développant une factorisation ou en substituant les racines.\par
\end{document}
